% !TEX TS-program = lualatex
% !TEX encoding = UTF-8

\documentclass[VesperaleOP.tex]{subfiles}

\ifcsname preamble@file\endcsname
  \setcounter{page}{\getpagerefnumber{M-vop11_tempus_post_trinitatem}}
\fi

\begin{document}

%%%% IN FESTO SS. TRINITATIS
\festum{}{In festo SS. Trinitatis}{I classis}{In festo SS. Trinitatis}{1}

\titulumrubrica{Ad I Vesperas}
\rubrica{Omnia ut in Psalterio, pag. \pageref{SVA1}.}

\rubrica{Post capitulum dicatur sequens Responsorium:}
\cantus{E1F7VR}{\rrann}
\rubrica{Ad Magnificat Ant. \normaltext{Peccáta}, pag. \pageref{SVAM}.}

\rubrica{Idem servatur in aliis Sabbatis usque ad Septuagesimam, nisi quod \rrrub hodie tantum dicatur.}
\oratio{Omnípotens sempitérne Deus, qui cæléstia simul et terréna moderáris:\psstar{} supplicatiónes pópuli tui cleménter exáudi;\pscross{}et pacem tuam nostris concéde tempóribus.\psstar{} Per.}

\titulumrubrica{Ad II Vesperas}
\rubrica{Omnia ut in Psalterio, pag. \pageref{DVA1}.}
\cantus{E2F1VAM}{Ad Magn.}
\rubrica{Oratio ut in I Vesperis.}
\lineamagna

%% ajout temporaire pour tester le rendu dans la table des matières
\festum{}{Antiphonæ propriæ in Dominicis}{A Dom. II post oct. Trinit. usque ad Adventum}{}{1}
\addcontentsline{toc}{section}{Antiphonæ in II Vesperis Dominicis post Pentecosten}

\end{document}